\documentclass[10pt]{article}
\usepackage[T1]{fontenc}
\usepackage[utf8]{inputenc}
\usepackage[french]{babel}
\usepackage{amsmath,amssymb,mathrsfs,mathpazo}
\usepackage{geometry}
\geometry{a4paper,margin=2cm}
\usepackage{graphicx,array,tabularx,tabularray,colortbl,multirow,multicol,booktabs,enumitem,pifont,ulem,fancybox,diagbox,cancel}
\usepackage{pstricks,pstricks-add,pst-plot,pst-tree,pst-node,pst-text,pst-func,pst-math,pst-bspline,pst-3dplot,pst-eucl,pst-solides3d}
\let\pscalculate\undefined
\usepackage{pst-bezier}
\usepackage{tikz,pgfplots,tkz-tab}
\usetikzlibrary{arrows,arrows.meta,patterns,calc,positioning,trees,decorations.markings,shapes.geometric,shapes.misc,shapes.arrows,intersections,decorations.pathmorphing,decorations.pathreplacing,backgrounds,shadows,babel}
\pgfplotsset{compat=1.18}
\usepgfplotslibrary{fillbetween}
\usepackage[np]{numprint}
\usepackage{listings,makecell,pas-tableur,twoopt}
\let\barre\undefined
\usepackage[pstricks]{tablvar}
\definecolor{darkgreen}{RGB}{0,100,0}
\providecommand{\red}{\color{red}}
\providecommand{\blue}{\color{blue}}
\providecommand{\green}{\color{green}}
\providecommand{\black}{\color{black}}
\providecommand{\white}{\color{white}}
\providecommand{\gray}{\color{gray}}
\providecommand{\cyan}{\color{cyan}}
\providecommand{\magenta}{\color{magenta}}
\providecommand{\yellow}{\color{yellow}}
\providecommand{\violet}{\color{violet}}
\providecommand{\orange}{\color{orange}}
\providecommand{\euro}{\texteuro}
\providecommand{\EUR}{\texteuro}
\providecommand{\ieme}{\textsuperscript{e}}
\providecommand{\degres}{\ensuremath{^\circ}}
\providecommand{\degree}{\ensuremath{^\circ}}
\providecommand{\vv}[1]{\overrightarrow{#1}}
\providecommand{\Coord}[1]{#1}
\providecommand{\decosix}{}
\providecommand{\starredbullet}{\ensuremath{\bullet}}
\pagestyle{empty}
\setlength{\parindent}{0pt}
\usepackage{adjustbox}\DeclareUnicodeCharacter{2236}{:}\DeclareUnicodeCharacter{2190}{\ensuremath{\leftarrow}}
\providecommand{\R}{\mathbb{R}}
\providecommand{\N}{\mathbb{N}}
\providecommand{\Z}{\mathbb{Z}}
\providecommand{\Q}{\mathbb{Q}}
\providecommand{\D}{\mathbb{D}}
\providecommand{\C}{\mathbb{C}}
\providecommand{\e}{{\mathrm{e}}}
\providecommand{\vect}[1]{\overrightarrow{#1}}
\providecommand{\vectt}[1]{\overrightarrow{\mathrm{#1}}}
\providecommand{\vvt}[1]{\overrightarrow{\mathrm{#1}}}
\providecommand{\np}[2][]{{#2~#1}}
\providecommand{\vv}[1]{{\overrightarrow{#1}}}
\renewcommand{\d}{{\mathrm{d}}}
\providecommand{\degres}{{^\circ}}
\providecommand{\degre}{{^\circ}}
\providecommand{\blue}{\color{blue}}
\providecommand{\red}{\color{red}}
\providecommand{\enskip}{\quad}
\providecommand{\llbracket}{\lbrack\!\lbrack}
\providecommand{\rrbracket}{\rbrack\!\rbrack}
\providecommand{\euro}{{\text{€}}}
\providecommand{\no}{{\text{n°}}}
\providecommand{\interval}[2]{{[#1;#2]}}
\providecommand{\intervFO}[2]{{[#1;#2[}}
\providecommand{\dotfill}{\ldots}

\providecommand{\ds}{\displaystyle}
\providecommand{\textding}[1]{#1}
\providecommand{\Oij}{\(\mathrm{O};\vec{\imath},\vec{\jmath}\)}
\providecommand{\Oijk}{\(\mathrm{O};\vec{\imath},\vec{\jmath},\vec{k}\)}
\providecommand{\Ouv}{\(\mathrm{O};\vec{u},\vec{v}\)}
\providecommand{\intervFF}[2]{\ensuremath{[#1;#2]}}
\renewcommand{\euro}{\ifmmode\text{\texteuro}\else\texteuro\fi}
\DeclareUnicodeCharacter{00B0}{\ifmmode^{\circ}\else\textdegree\fi}
\DeclareUnicodeCharacter{20AC}{\ifmmode\text{\texteuro}\else\texteuro\fi}
% Énoncés : Éducation nationale. Transcription et corrigés : APMEP, Denis Vergès et contributeurs. Découpage et classement : CoopMaths, Sébastien Lozano, Éric Elter et contributeurs. Adaptation de la présentation : Maths Explorer.
% CC BY-SA : https://coopmaths.fr/www/about/
% https://coopmaths.fr/alea/static/e3c/2020/tex/e3c_2020_00_sujet9_3_cor.tex
\begin{document}
\noindent\textit{APMEP / CoopMaths. CC BY-SA. Présentation adaptée par Maths Explorer.}\par\medskip
\section*{Exercice 3 (5 points)}	
	\subsection*{1. Déterminer par lecture graphique les coordonnées des points $A$, $B$, $C$, $K$ et $L$.}
	
	\[
	A(0; 35), B(24; 35), C(24; 0), K(12; 35), L(24; 7)
	\]
	
	\subsection*{2. Un visiteur affirme : « Moins de 70 \% de la surface de la place est éclairée ». Cette affirmation est-elle exacte ?}
	
	\[
	\begin{array}{l}
		\text{Aire éclairée} = \text{Aire}(OCBA) - \text{Aire}(OCL) - \text{Aire}(OKA) \\
		= 24 \times 35 - \dfrac{1}{2} \times 24 \times 7 - \dfrac{1}{2} \times 12 \times 35 \\
		= 24 \times 35 - 12 \times 7 - 6 \times 35 \\
		= 840 - 84 - 210 = 840 - 294 = 456 \text{ u.a.} \\
		\\
		\dfrac{456}{840} = \dfrac{19}{35} \approx 0,54 \text{ soit environ 54\% de la place est éclairée.}
	\end{array}
	\]
L’affirmation est exacte.
	
	\subsection*{3. a. Donner les coordonnées des vecteurs $\overrightarrow{OK}$ et $\overrightarrow{OL}$.}
	
	\[
	\overrightarrow{OK} \left( \begin{array}{c} 12 \\ 35 \end{array} \right) \quad \text{et} \quad \overrightarrow{OL} \left( \begin{array}{c} 24 \\ 7 \end{array} \right)
	\]
	
	\subsection*{b. Montrer que le produit scalaire $\overrightarrow{OK} \cdot \overrightarrow{OL}$ est égal à 533.}
	
	\[
	\overrightarrow{OK} \cdot \overrightarrow{OL} = 12 \times 24 + 35 \times 7 = 288 + 245 = 533
	\]
	
	\subsection*{c. En déduire la mesure, arrondie au degré, de l’angle $\widehat{KOL}$.}
On sait que l’on a aussi :
	\[
	\overrightarrow{OK} \cdot \overrightarrow{OL} = OK \times OL \times \cos(\overrightarrow{OK}; \overrightarrow{OL})
	\]
Avec :
	\[
	OK = \sqrt{12^2 + 35^2} = \sqrt{144 + 1225} = \sqrt{1369} = 37
	\]
	\[
	OL = \sqrt{24^2 + 7^2} = \sqrt{576 + 49} = \sqrt{625} = 25
	\]
L’égalité devient :
	\[
	533 = 37 \times 25 \times \cos(\overrightarrow{OK}; \overrightarrow{OL}) \quad \text{d’où} \quad \cos(\overrightarrow{OK}; \overrightarrow{OL}) = \dfrac{533}{37 \times 25} = \dfrac{533}{925}
	\]
La calculatrice donne $\widehat{KOL} \approx 54,81$ soit $55^\circ$ au degré près.

\end{document}
